\documentclass[11pt]{article}
\usepackage[utf8]{inputenc}
\usepackage[T1]{fontenc}
\usepackage[margin=1in]{geometry}
\usepackage{parskip}

\title{Qihang Han Presentation Script}
\author{}
\date{}

\begin{document}

\maketitle

\section*{Qihang Han --- On-Chain Vault, Live Demo, and Closing}
\textbf{Target speaking time: 3:55}

\subsection*{Slide: On-Chain Deep Dive}
\textbf{Time: 8:30--9:05}

The Solidity side is centered around two contracts.

\texttt{MockStablecoin} gives us a local demo asset, so we can show mint, approve, deposit, and withdraw flows without relying on a live token.

\texttt{DeltaNeutralVault} is the main contract. It handles deposits, withdrawals, internal share accounting, strategy-state updates, NAV synchronization, PnL updates, pause control, operator management, and events.

The key point is that we are not executing the full strategy on-chain. The vault is the accounting and state-management layer for an off-chain strategy engine.

\subsection*{Slide: Boundary Handoff}
\textbf{Time: 9:05--9:35}

The next question is what actually crosses from off-chain analytics into on-chain state.

We keep the payload small. The contract does not receive the full dataset, model, or backtest. It receives only the fields needed for state:

strategy state, reported NAV, PnL delta, signal hash, metadata hash, and report hash.

The hashes anchor the on-chain update to off-chain evidence. The vault stores compact state, while the detailed reasoning stays in generated artifacts and reports.

This keeps the design lightweight, makes the trust boundary explicit, and gives the blockchain a meaningful role without forcing unsuitable computation onto the EVM.

\subsection*{Slide: Implementation Status}
\textbf{Time: 9:35--10:00}

In terms of implementation, the main demo layers are already built.

Off-chain, we already have data, features, labels, models, signals, backtest, evaluation, reporting, demo, and integration modules.

On-chain, we already have the mock stablecoin, the delta-neutral vault, update functions, events, pause logic, ownership mechanics, and operator control.

Around that, we also have local deployment and update scripts, Foundry tests, and the frontend bridge that packages everything into the presentation page and dashboard.

\subsection*{Slide: Blockchain Evidence}
\textbf{Time: 10:00--10:25}

This slide answers a simple question: how do we know the chain side is more than a diagram?

The answer is that the vault exposes concrete capabilities: deposit and withdraw, owner and operator separation, strategy-state updates, NAV and PnL synchronization, pause control, and event logging.

These are not just labels on a slide. A deposit mints shares, a NAV update changes reported assets, and a strategy update records state plus hashes. Events provide the audit trail.

Now I will use the live vault console to show the state transitions directly.

\subsection*{Slide: Live Vault Console}
\textbf{Time: 10:25--11:40}

This is a local presentation simulator of the vault workflow. It mirrors our Solidity vault logic, but it is stable for classroom use because it does not depend on a wallet or live RPC connection.

I will move through four actions quickly.

\textbf{Click: Deploy local vault}

This deploys \texttt{MockStablecoin} and \texttt{DeltaNeutralVault}, and we can see the vault address and event record appear.

\textbf{Click: Deposit 10,000 mUSDC}

This represents user funding. Wallet balance decreases, vault cash increases, reported NAV increases, and shares are minted.

\textbf{Click: Sync strategy state}

This pushes the selected strategy state from off-chain analytics into the contract-facing state.

\textbf{Click: Sync NAV / PnL}

This synchronizes the valuation result. Reported NAV, cumulative PnL, and the report hash all update.

\textbf{Click: Withdraw 2,500 mUSDC}

Finally, withdrawal burns shares and returns assets to the user. The main point of this demo is accountability: user-facing state, vault accounting, and event-style logs all change step by step.

\subsection*{Slide: Chain Demo Flow / Vault Demo Story}
\textbf{Time: 11:40--12:00}

This slide summarizes the same process at the system level.

The flow is deployment, user funding and approval, deposit, strategy-state update, NAV or PnL synchronization, and withdrawal.

The vault stores the parts that should be transparent to users: shares, reported NAV, strategy state, timestamps, and hash references to off-chain artifacts.

So the frontend is not just showing charts. It connects research output to a contract-facing state machine.

\subsection*{Slide: Caveats}
\textbf{Time: 12:00--12:15}

There are clear limitations. This is a course-project prototype, not an audited DeFi protocol. The vault still uses a trusted owner or operator path, and execution remains off-chain.

These limitations are intentional. The goal is to demonstrate a credible hybrid architecture, not to overclaim production readiness.

\subsection*{Slide: Next Steps / Closing}
\textbf{Time: 12:15--12:25}

To conclude, our project contributes four things.

First, an end-to-end off-chain research pipeline for funding-rate arbitrage. Second, a cost-aware strategy evaluation workflow. Third, a Solidity vault that represents pooled capital, shares, strategy state, NAV, and PnL. Fourth, a frontend demo that connects the research system to the vault interface.

Future work would include signed or oracle-mediated updates, richer vault lifecycle controls, multi-asset support, and a cleaner operator dashboard.

Thank you. We are happy to take questions.

\end{document}
